\section{Core Architecture}

\SCS{} uses a layered enforcement architecture.

\begin{lstlisting}
External transaction
    |
    v
Contract Dispatcher
    |
    v
Method Policy Registry
    |
    v
Enforcement Module
    |
    v
Live Capability Registry
    |
    v
Authorized Execution Frame
    |
    v
Private Method Implementation
    |
    v
Guarded Storage Kernel
    |
    v
Atomic Transaction Engine
    |
    v
Committed State + Events
\end{lstlisting}

The runtime has five main enforcement layers:

\begin{description}[leftmargin=2.8cm]
  \item[Dispatcher] controls which methods are callable.
  \item[Method policy registry] declares required authority and allowed effects.
  \item[Enforcement module] checks current caller against live registry state.
  \item[Guarded storage kernel] rejects raw mutation and writes outside method scope.
  \item[Atomic transaction engine] commits all registry, state, and event effects together.
\end{description}

\subsection{Ordinary call flow}

A normal method call follows this shape:

\begin{lstlisting}
(call-contract! tx TokenA transferFrom Alice Bob USDC 10)
\end{lstlisting}

The caller does not pass a capability token. Instead, the runtime binds the transaction context, finds the method policy for \code{TokenA.transferFrom}, derives the required allowance from the method arguments, and checks the live capability registry.

\subsection{State transition boundary}

The private implementation may request state reads and writes through runtime primitives. The storage kernel accepts a write only if:

\begin{itemize}
  \item an authorized execution frame is active;
  \item the write targets the current contract instance;
  \item the key matches the active method's write scope;
  \item the value satisfies schema and kernel invariants;
  \item the transaction engine can stage the write deterministically.
\end{itemize}
